\subsection{Packing the wire slots}\label{sec:compressor-slots}

The lane coordinates let us replicate each base arc in four endpoint
configurations. We mark each endpoint thick or thin. Thick positions
will merge at a fixed site, whereas thin positions can move to either
end of an assigned bundle. The assignment must move every occurrence
of a given old position together, even when several motifs meet there.

\begin{lemma}[Packing and placement of the wire slots]\label{lem:compressor-slots}
For every power-of-two base degree $d$, set $K=2^{2d+1}d$ and $Q=3K/4$.
There is a packing of the wire slots, using a number of lane coordinates
depending only on $d$, with the following properties. The two layouts
refer to the six-party motif of Figure~\ref{fig:motif-layout}.
\begin{enumerate}
\item Packing the thick positions produces $P/K$ new positions. Each
thick--thick bundle contains $K/4$ wires. Every other wire is assigned as
a $CD$ demand to the bundle of its thin $C$ endpoint, with $Q$ per bundle.
\item In either layout the packed endpoints $A,B$ and the image $D^*=\pi(D)$
of the other endpoint have distinct bases in each motif. Each of these three
roles occurs exactly $dQ$ times per new position.
\item Translation of the remaining lane preserves the assignments, placements
and unweighted slots. Slot numbering, including the padding
to any admissible output degree, can be determined by bounded-size lists
per new position after a greedy scan of the base.
\end{enumerate}
\end{lemma}

\begin{proof}
Color the base greedily with $2d+1$ colors. Split the lane as $(U,h,j,g)$,
where $U$ supplies one bit per color, $h$ indexes the $d$ outgoing slots,
$j$ is a tag in a binary vector space of order $J$, and $g$ is the remaining
lane. We determine $J$ below. Call a position thick if its base-color bit
is $1$, and thin otherwise. Put $K=2^{2d+1}d$, the number of pairs $(U,h)$.
Along a base arc the endpoint colors are distinct. With $j,g$ fixed at one
end, the two tests use distinct free bits of $U$ up to translations. They
are independent fair tests, also independently of $h$ at either end.
Consequently each of the four endpoint types occurs $K/4$ times.

For given $x,j,g$ pack all thick positions there into a single new one, with
base $(x,j)$ and lane $g$. There are $P'=P/K$ new positions. Bundle the
$K/4$ thick--thick wires of each base arc for fixed $j,g$ at its source by
concatenating them in a common order at both endpoints. The bundle runs from
a packed position $A$ to a packed position $B$. For every wire outside these
bundles take a thin end for $C$ and the other end for $D$. For thin--thin
wires take their old direction $C\to D$.

Give each base vertex $x$ an offset $\phi_x$. Assign each thin position
$(x,U,h,j,g)$ to the bundle with source $(x,j+\phi_x,g)$ on the outgoing
base arc of slot $h$. Every wire on which this position plays $C$ is assigned
to that bundle. To count these wires, fix the base, tag, remaining lane and
slot of its possible thin sources. The $d$ originally outgoing incidences
contribute $d(K/(2d))=K/2$ wires. The $d$ originally incoming incidences
from thick positions contribute $d(K/(4d))=K/4$. Thus each bundle receives
$Q=3K/4$ wires of type $CD$.

We choose the tag offsets so that $A,B,D^*$ always have distinct new bases.
The old base components of $A,B$ differ. For $C$ at old base $x$ with tag
$j$, and $D$ at $y$ on their old base arc $e$ in either direction, write
\[
          L_x=\{0\}\cup\{\tau_f^{\mathrm{tag}}: f\text{ leaves }x\}.
\]
The new tags of $A,B$ lie in $j+\phi_x+L_x$. A thick $D$ has tag
$j+\tau_e^{\mathrm{tag}}$, and the possible tags of a thin $D^*$ belong to
$j+\tau_e^{\mathrm{tag}}+\phi_y+L_y$. It suffices to impose, at each
incidence,
\[
      \phi_x\notin\tau_e^{\mathrm{tag}}+L_x,
      \qquad \phi_x+\phi_y\notin\tau_e^{\mathrm{tag}}+L_x+L_y.
\]
Greedy choice in base order works for any power of two
$J>2d((d+1)+(d+1)^2)$, since at most this many values are forbidden by
the first conditions and the already assigned neighbors in the second.

\smallskip
At every new position $d$ bundles start and $d$ bundles end, so each of
the roles $A,B$ occurs $dQ$ times. The thick $D$ endpoints at this position
contribute $2dK/4=dK/2$ occurrences of $D^*$, counting mixed arcs in both
old directions. The others are destinations of thin--thin arcs. Before
placement, for fixed base, tag and remaining lane, these contribute $dK/4$
occurrences, equally distributed over destination slot labels $h$. In
layout~1 the tag offset permutes the tags. In layout~2 each position receives
through $d$ incoming base arcs, each with a uniquely determined source tag
and remaining lane, and each carrying $K/4$ such occurrences from its
assigned slot. In both cases the total number of $D^*$ roles is
$dK/2+dK/4=dQ$.

Packing, bundle assignment, and the two choices of image depend on the
remaining lane only through the given arc translations. The common choice
of layout therefore preserves this symmetry. For each surviving occurrence,
label its parallel slots by the original $CD$ arc identifying the motif, the
surviving type, and the subdivision index. When a common deficit remains
after the degree calculation in Subsection~\ref{subsec:shrink}, add zero-width
arcs to the next tag cyclically at the same old base and remaining lane.
Number incoming and outgoing slots separately by sorting their identifiers,
expressing all remaining lanes relative to the vertex being numbered. The
orders transport from lane zero to every lane. Widths and term-dependent map
entries need never enter these identifiers.

Greedy coloring and tag selection each use a single scan of the indexed
base and bounded neighbor lists. After these scans the assignments can be
formed by listing the actual arcs by bundle, and the regularized slots by
their endpoints. The list sizes depend only on the fixed degrees and lane
coordinates consumed in this step. The sorted identifiers give the two
corresponding slot numbers and their displacement, as required for the
next level of the construction.
\end{proof}
